\section{Failure analysis}
\label{sec:failure}

La consegna chiede almeno un caso di fallimento: ne vengono riportati \textbf{tre},
ottenuti portando all'estremo un iperparametro alla volta su CartPole con
Q\hbox{-}Learning (Fig.~\ref{fig:cp-fail}). Ciascuno \`e scelto perch\'e rompe un
meccanismo \emph{diverso}: il passo di aggiornamento, l'esplorazione e il peso del futuro. 
In tutti e tre i casi il codice gira senza errori, giugneendo alla conclusione che 
la correttezza dell'implementazione non garantisce l'apprendimento.

\paragraph{1. $\alpha = 1.0$ --- instabilit\`a (mancata convergenza).}
L'aggiornamento \emph{sostituisce} completamente la stima ad ogni passo
($q(S_t,A_t) \leftarrow R_{t+1} + \gamma \max_{a'} q(S_{t+1},a')$) e l'agente si blocca a $\approx 65$
per l'intera durata dei 6\,000 episodi.

La lettura pi\`u importante riguarda per\`o un aspetto diverso: il fallimento \`e specifico
del valore \emph{estremo}, non dell'``$\alpha$ alto'' in generale. Nella stessa figura
$\alpha = 0.5$ non solo non fallisce, ma risulta la curva \emph{migliore}
($\approx 210$, superiore ad $\alpha=0.1$). Se avessi testato soltanto $\{0.1, 1.0\}$
avrei concluso, erroneamente, che ``$\alpha$ grande destabilizza'': la transizione non
\`e graduale e la soglia di instabilit\`a si colloca tra $0.5$ e $1.0$.

\paragraph{2. $\varepsilon = 0$ --- comportamento subottimale.}
L'agente opera in modo greedy fin dal primo istante su una Q-table inizializzata a zero:
si fissa sulla prima azione che sembra funzionare senza mai provare le alternative. La
curva resta piatta a $\approx 40$, mentre con $\varepsilon = 0.1$ si raggiunge
$\approx 170$. Degno di nota il caso intermedio $\varepsilon = 0.01$, che non fallisce
del tutto ma resta bloccato per $\approx 2\,500$ episodi prima di sbloccarsi e arrivare
a $\approx 100$. Un'esplorazione insufficiente non impedisce l'apprendimento in senso
assoluto, ma lo rende estremamente inefficiente (S\&B \S2.1).

\paragraph{3. $\gamma = 0$ --- miopia.}
Con $\gamma = 0$ il target si
riduce al reward immediato; tuttavia, su CartPole il reward immediato \`e $+1$ ad
\emph{ogni} passo non terminale, \emph{qualunque} azione venga scelta. Tutte le coppie
stato-azione collassano pertanto sullo stesso valore, la policy greedy non ha modo di
preferire un'azione all'altra e degenera in una policy casuale: la curva resta piatta a
$\approx 22$, vale a dire \emph{al di sotto} del punto di partenza delle altre
configurazioni.

Questo caso mostra concretamente \emph{a cosa serve} il discount factor: non
semplicemente a ``pesare meno il futuro'', ma a creare il canale attraverso cui
l'informazione sulla caduta futura risale fino alle azioni che l'hanno causata. Senza
$\gamma > 0$ quel canale non esiste, e nessun numero di episodi pu\`o rimediare.

\begin{figure}[htbp]
\centering
\begin{subfigure}{0.31\linewidth}
  \includegraphics[width=\linewidth]{cartpole_failure_high_alpha.png}
  \caption{$\alpha=1.0$: instabilit\`a.}
\end{subfigure}\hfill
\begin{subfigure}{0.31\linewidth}
  \includegraphics[width=\linewidth]{cartpole_failure_no_exploration.png}
  \caption{$\varepsilon=0$: subottimale.}
\end{subfigure}\hfill
\begin{subfigure}{0.31\linewidth}
  \includegraphics[width=\linewidth]{cartpole_failure_no_discount.png}
  \caption{$\gamma=0$: miopia.}
\end{subfigure}
\caption{CartPole, i tre casi di fallimento (Q\hbox{-}Learning, media su 3 seed).
(a) $\alpha=1.0$ (verde) si blocca a $\approx 65$, mentre $\alpha=0.5$ (arancio) \`e la
curva migliore: la soglia di instabilit\`a si colloca tra i due valori, non
genericamente ``in alto''. (b) $\varepsilon=0$ (blu) resta piatto a $\approx 40$, mentre
$\varepsilon=0.01$ (arancio) apprende con $\approx 2\,500$ episodi di ritardo.
(c) $\gamma=0$ (blu) \`e piatto a $\approx 22$: con reward $+1$ ad ogni passo tutte le
azioni risultano equivalenti e la policy greedy \`e indistinguibile da una casuale.}
\label{fig:cp-fail}
\end{figure}

